%% kapa: exam=1 round=28 year=2017 period=1 form=A subject=real_estate subject_ko=부동산학원론 seq=37 number=117
%% source: Q-Net 감정평가사 기출문제 (cst00309) artlSeq=5199603 file=extracted/1008/5199603/2017년 제28회 감정평가사 1차시험 문제지/1교시 A형.pdf
%% conversion: kordoc 4.24.51 -> tools/gen1.py (auto); review=proofread-sonnet; flags=none
\begin{question}{117}
택지개발방식 중 환지방식에 관한 설명으로 옳지 않은 것을 모두 고른 것은?
\begin{framedtext}
ㄱ. 사업자로서는 상대적으로 사업시행이 간단하고 용이하다. \par
ㄴ. 개발이익은 토지소유자, 사업자 등이 향유한다. \par
ㄷ. 사업자의 초기 사업비 부담이 크고, 토지소유자의 저항이 심할 수 있다. \par
ㄹ. 감보된 토지는 새로이 필요로 하는 공공시설 용지로 사용되고, 나머지 체비지는 환지한다. \par
ㅁ. 환지의 형평성을 기하기 위해 사업시행기간이 장기화 될 수 있다. \par
ㅂ. 혼용방식은 수용 또는 사용방식과 환지방식을 혼용하여 시행하는 방식이다.
\end{framedtext}
\begin{choices}
  \item ㄱ, ㄴ, ㄷ
  \item ㄱ, ㄷ, ㄹ
  \item ㄱ, ㄹ, ㅁ
  \item ㄴ, ㅁ, ㅂ
  \item ㄹ, ㅁ, ㅂ
\end{choices}
\end{question}
\answer{2}
